\documentclass[10pt,a4paper]{amsart}
\usepackage{amsmath,amssymb,amsthm,graphicx}
\usepackage{verbatim}
\newenvironment{reference}{\comment}{\endcomment}
%\newenvironment{reference}{}{}
\newenvironment{slogan}{\comment}{\endcomment}
\newenvironment{history}{\comment}{\endcomment}
% For commutative diagrams we use Xy-pic
\usepackage[all]{xy}

% We use 2cell for 2-commutative diagrams.
\xyoption{2cell}
% Theorem environments.
%
\theoremstyle{plain}
\newtheorem{theorem}[subsection]{Theorem}
\newtheorem{proposition}[subsection]{Proposition}
\newtheorem{lemma}[subsection]{Lemma}

\theoremstyle{definition}
\newtheorem{definition}[subsection]{Definition}
\newtheorem{example}[subsection]{Example}
\newtheorem{exercise}[subsection]{Exercise}
\newtheorem{situation}[subsection]{Situation}

\theoremstyle{remark}
\newtheorem{remark}[subsection]{Remark}
\newtheorem{remarks}[subsection]{Remarks}

\numberwithin{equation}{subsection}

% Macros
%
\def\lim{\mathop{\mathrm{lim}}\nolimits}
\def\colim{\mathop{\mathrm{colim}}\nolimits}
\def\Spec{\mathop{\mathrm{Spec}}}
\def\Hom{\mathop{\mathrm{Hom}}\nolimits}
\def\Ext{\mathop{\mathrm{Ext}}\nolimits}
\def\SheafHom{\mathop{\mathcal{H}\!\mathit{om}}\nolimits}
\def\SheafExt{\mathop{\mathcal{E}\!\mathit{xt}}\nolimits}
\def\Sch{\mathit{Sch}}
\def\Mor{\mathop{\mathrm{Mor}}\nolimits}
\def\Ob{\mathop{\mathrm{Ob}}\nolimits}
\def\Sh{\mathop{\mathit{Sh}}\nolimits}
\def\NL{\mathop{N\!L}\nolimits}
\def\CH{\mathop{\mathrm{CH}}\nolimits}
\def\proetale{{pro\text{-}\acute{e}tale}}
\def\etale{{\acute{e}tale}}
\def\QCoh{\mathit{QCoh}}
\def\Ker{\mathop{\mathrm{Ker}}}
\def\Im{\mathop{\mathrm{Im}}}
\def\Coker{\mathop{\mathrm{Coker}}}
\def\Coim{\mathop{\mathrm{Coim}}}

% Boxtimes
%
\DeclareMathSymbol{\boxtimes}{\mathbin}{AMSa}{"02}

%
% Macros for moduli stacks/spaces
%
\def\QCohstack{\mathcal{QC}\!\mathit{oh}}
\def\Cohstack{\mathcal{C}\!\mathit{oh}}
\def\Spacesstack{\mathcal{S}\!\mathit{paces}}
\def\Quotfunctor{\mathrm{Quot}}
\def\Hilbfunctor{\mathrm{Hilb}}
\def\Curvesstack{\mathcal{C}\!\mathit{urves}}
\def\Polarizedstack{\mathcal{P}\!\mathit{olarized}}
\def\Complexesstack{\mathcal{C}\!\mathit{omplexes}}
% \Pic is the operator that assigns to X its picard group, usage \Pic(X)
% \Picardstack_{X/B} denotes the Picard stack of X over B
% \Picardfunctor_{X/B} denotes the Picard functor of X over B
\def\Pic{\mathop{\mathrm{Pic}}\nolimits}
\def\Picardstack{\mathcal{P}\!\mathit{ic}}
\def\Picardfunctor{\mathrm{Pic}}
\def\Deformationcategory{\mathcal{D}\!\mathit{ef}}
\usepackage{cleveref}
\crefname{originalreference}{original reference}{original references}
\Crefname{originalreference}{Original reference}{Original references}
\makeatletter
\newlabel{divisors-lemma-bourbaki}{{divisors-lemma-bourbaki}{source}{}{sourcepointer-1}{}}
\newlabel{divisors-lemma-bourbaki@cref}{{[originalreference][0][]divisors-lemma-bourbaki}{[1][0][]source}{}{}{}}
\makeatother
\crefname{originalreference}{original reference}{original references}
\Crefname{originalreference}{Original reference}{Original references}
\makeatletter
\newlabel{dualizing-lemma-divide-by-torsion}{{dualizing-lemma-divide-by-torsion}{source}{}{sourcepointer-2}{}}
\newlabel{dualizing-lemma-divide-by-torsion@cref}{{[originalreference][0][]dualizing-lemma-divide-by-torsion}{[1][0][]source}{}{}{}}
\makeatother
\crefname{originalsection}{original section}{original sections}
\Crefname{originalsection}{Original section}{Original sections}
\makeatletter
\newlabel{section-supports}{{section-supports}{source}{}{sourcepointer-3}{}}
\newlabel{section-supports@cref}{{[originalsection][0][]section-supports}{[1][0][]source}{}{}{}}
\makeatother
\title{The universal quotient of a coherent module away from the points of a specialization-stable subset, with depth at least one there}
\author{The Stacks Project Authors}
\date{stacks-project:local-cohomology.tex:local-snapshot-20261012; GNU FDL 1.2 or later}
\begin{document}
\maketitle
\noindent\textit{Editorial note.} All delivered mathematical text is the author's own, copied from the pinned Stacks Project chapter source sources/stacks-local-cohomology/pinned-original-source.tex; no statement or proof step was rewritten, reordered or removed.\par\medskip
\noindent\textit{Editorial note.} Preamble provenance: the Stacks Project keeps one single shared preamble for every chapter; this card takes it directly from the staged sources/stacks-local-cohomology/preamble.tex as three verbatim spans (lines 8-12, 14-18 and 154-217), so no derived file was written and nothing was added or redefined. The wrapper class and the base packages are supplied by the assembler.\par\medskip
\noindent\textit{Editorial note.} Bibliography provenance: neither the counted statement nor its proof contains a citation (this lane rescans the whole span before writing the card), so the card carries no bibliography and no entry was reconstructed.\par\medskip
\noindent\textit{Editorial note.} Selection provenance: this card was opened under the file-based control contract (CONTROL.md), with the mechanical pre-card gate run on its statement and proof interval before the card was written.\par\medskip
\noindent\textit{Allowed standard prerequisites (editorial).}\begin{itemize}
\item The subsheaf of sections whose support is contained in a specialization-stable subset, and its description as the filtered colimit of the sections with support in the closed subsets contained in it. The author's section on supports of modules and its local description (source line 3325), used together with a lemma of the chapter on dualizing complexes. Identifies the kernel of the required quotient and makes the depth condition a condition on associated primes.
\item Flat base change for such quotients, in the form recorded in the chapter on divisors. Invoked inside the counted proof at source line 3315-3316 as a cross-chapter pointer. Gives the compatibility of the construction with flat morphisms once uniqueness is known.
\end{itemize}\medskip
\ifdefined\hypertarget\hypertarget{sourcepointer-1}{}\fi
\noindent\textit{Source reference.} Original reference divisors-lemma-bourbaki: sources/stacks-local-cohomology/pinned-original-source.tex, source lines 3316--3316. Navigation pointer only: the counted text refers to the Stacks item divisors-lemma-bourbaki, which belongs to another chapter of the project and is not part of this card. The pointer carries the target's own label as its number because this chapter's staged snapshot does not contain the other chapter's tag table; its content is not reproduced here.\par\medskip
\ifdefined\hypertarget\hypertarget{sourcepointer-2}{}\fi
\noindent\textit{Source reference.} Original reference dualizing-lemma-divide-by-torsion: sources/stacks-local-cohomology/pinned-original-source.tex, source lines 1513--1513. Navigation pointer only: the counted text refers to the Stacks item dualizing-lemma-divide-by-torsion, which belongs to another chapter of the project and is not part of this card. The pointer carries the target's own label as its number because this chapter's staged snapshot does not contain the other chapter's tag table; its content is not reproduced here.\par\medskip
\ifdefined\hypertarget\hypertarget{sourcepointer-3}{}\fi
\noindent\textit{Source reference.} Original section section-supports: sources/stacks-local-cohomology/pinned-original-source.tex, source lines 3325--3325. Navigation pointer only: the counted text refers to the Stacks item section-supports, which belongs to another chapter of the project and is not part of this card. The pointer carries the target's own label as its number because this chapter's staged snapshot does not contain the other chapter's tag table; its content is not reproduced here.\par\medskip
\begin{lemma}
\label{lemma-get-depth-1-along-Z}
Let $X$ be a Noetherian scheme. Let $T \subset X$ be a subset
stable under specialization. Let $\mathcal{F}$ be a coherent
$\mathcal{O}_X$-module. Then there is a unique map
$\mathcal{F} \to \mathcal{F}'$ of coherent $\mathcal{O}_X$-modules
such that
\begin{enumerate}
\item $\mathcal{F} \to \mathcal{F}'$ is surjective,
\item $\mathcal{F}_x \to \mathcal{F}'_x$ is an isomorphism for $x \not \in T$,
\item $\text{depth}_{\mathcal{O}_{X, x}}(\mathcal{F}'_x) \geq 1$ for $x \in T$.
\end{enumerate}
If $f : Y \to X$ is a flat morphism with $Y$ Noetherian, then
$f^*\mathcal{F} \to f^*\mathcal{F}'$ is the corresponding
quotient for $f^{-1}(T) \subset Y$ and $f^*\mathcal{F}$.
\end{lemma}

\begin{proof}
Condition (3) just means that $\text{Ass}(\mathcal{F}') \cap T = \emptyset$.
Thus $\mathcal{F} \to \mathcal{F}'$ is the quotient of $\mathcal{F}$
by the subsheaf of sections whose support is contained in $T$.
This proves uniqueness. The statement on pullbacks follows from
Divisors, Lemma \ref{divisors-lemma-bourbaki}
and the uniqueness.

\medskip\noindent
Existence of $\mathcal{F} \to \mathcal{F}'$.
By the uniqueness it suffices to prove the
existence and uniqueness locally on $X$; small detail omitted.
Thus we may assume $X = \Spec(A)$ is affine and $\mathcal{F}$
is the coherent module associated to the finite $A$-module $M$.
Set $M' = M / H^0_T(M)$ with $H^0_T(M)$ as in Section \ref{section-supports}.
Then $M_\mathfrak p = M'_\mathfrak p$ for $\mathfrak p \not \in T$
which proves (1). On the other hand, we have
$H^0_T(M) = \colim H^0_Z(M)$ where $Z$ runs over the closed
subsets of $X$ contained in $T$. Thus by
Dualizing Complexes, Lemmas \ref{dualizing-lemma-divide-by-torsion}
we have $H^0_T(M') = 0$, i.e., no associated prime
of $M'$ is in $T$. Therefore $\text{depth}(M'_\mathfrak p) \geq 1$
for $\mathfrak p \in T$.
\end{proof}

\end{document}
